\documentclass{article}
\usepackage{amsmath,amssymb}
\usepackage{xurl}
\usepackage[hidelinks]{hyperref}
% Shared commands for notes in docs/paper-gaps/.

\DeclareUnicodeCharacter{2080}{\ifmmode{}_{0}\else\textsubscript{0}\fi}
\DeclareUnicodeCharacter{2081}{\ifmmode{}_{1}\else\textsubscript{1}\fi}
\DeclareUnicodeCharacter{2082}{\ifmmode{}_{2}\else\textsubscript{2}\fi}
\DeclareUnicodeCharacter{2099}{\ifmmode{}_{n}\else\textsubscript{n}\fi}

\newcommand{\Lean}{\textsc{Lean}}
\ExplSyntaxOn
\NewDocumentCommand{\leanid}{m}{
  \ifmmode
    \textsf{\detokenize{#1}}
  \else
    \group_begin:
    \urlstyle{sf}
    \tl_set:Nn \l_tmpa_tl {#1}
    \tl_replace_all:Nnn \l_tmpa_tl {\_} {_}
    \exp_args:NV \path \l_tmpa_tl
    \group_end:
  \fi
}
\ExplSyntaxOff
\providecommand{\tr}{\operatorname{tr}}
\newcommand{\tnleanIssueBase}{https://github.com/LionSR/TNLean/issues/}
\newcommand{\tnleanPrBase}{https://github.com/LionSR/TNLean/pull/}
\newcommand{\tnleanDocBase}{https://sirui-lu.com/TNLean/docs/}
\newcommand{\ghissue}[1]{\href{\tnleanIssueBase#1}{GitHub issue \##1}}
\newcommand{\ghpr}[1]{\href{\tnleanPrBase#1}{PR \##1}}
\newcommand{\doclink}[1]{\href{\tnleanDocBase#1.html}{\path{#1}}}

% Dirac notation and the identity operator, as in blueprint/src/macros/common.tex.
% \providecommand keeps note-local definitions authoritative.
\usepackage{dsfont}
\providecommand{\ket}[1]{|#1\rangle}
\providecommand{\bra}[1]{\langle #1|}
\providecommand{\braket}[2]{\langle #1 | #2 \rangle}
\providecommand{\ketbra}[2]{|#1\rangle\!\langle #2|}
\providecommand{\Id}{\mathds{1}}
\providecommand{\one}{\mathds{1}}
\providecommand{\C}{\mathbb{C}}
\providecommand{\MN}[1]{M_{#1}(\C)}
\providecommand{\MD}{M_D(\C)}

% External referencing for paper sources.  The build helper writes sanitized
% auxiliary files under build/paper-aux/ and excludes duplicated source labels.
\usepackage{xr}

% Import a generated paper auxiliary file with a caller-supplied label prefix.
% The candidate paths support builds from docs/paper-gaps/, the repository root,
% and build/paper-aux/ itself.
\newcommand{\importpaperaux}[2]{%
  \IfFileExists{../../build/paper-aux/#2.aux}{%
    \externaldocument[#1]{../../build/paper-aux/#2}%
  }{%
    \IfFileExists{build/paper-aux/#2.aux}{%
      \externaldocument[#1]{build/paper-aux/#2}%
    }{%
      \IfFileExists{#2.aux}{%
        \externaldocument[#1]{#2}%
      }{}%
    }%
  }%
}

% General external reference macros
\newcommand{\paperref}[2]{\ref{#1-#2}}
\newcommand{\papereqref}[2]{(\ref{#1-#2})}

% CPSV16 source (1606.00608 / MPDO-22-12-17-2.tex) external document and shortcuts
\importpaperaux{cpsv16-}{1606.00608/MPDO-22-12-17-2-tnlean}
\newcommand{\cpsvref}[1]{\paperref{cpsv16}{#1}}
\newcommand{\cpsveqref}[1]{\papereqref{cpsv16}{#1}}

% Verdict marker: \gapnote{<kind>}{<status>}. Kind is the policy
% classification (clarification, local-correction, scope-restriction,
% unfaithful, false-source, open-gap); status is open, wip (elimination
% underway), resolved, or historical. Machine-read by the site index;
% prints a verdict line.
\newcommand{\gapnote}[2]{\par\noindent\textbf{Verdict:} #1 (#2).\par}

\title{Two-Dimensional PEPS Examples: Tori of Width Two}
\author{}
\date{2026-09-26}
\begin{document}
\maketitle
\gapnote{scope-restriction}{open}

\section{What the sources print}
Appendix~A of \cite{Cirac2021Matrix}, under ``Two dimensions: PEPS'', gives
single-site tensors with four virtual indices ordered top, right, down and
left, to be placed at every site of a square lattice.\footnote{Source
    traceability: \path{Papers/2011.12127/TN-Review-main.tex}, line 2415.}
The cluster state is written with the tensor
$A=\ket{0}(00{+}{+}|+\ket{1}(11{-}{-}|$ on the square lattice, with no
restriction on its size.\footnote{Source traceability:
    \path{Papers/2011.12127/TN-Review-main.tex}, lines 2424--2429.}
The CZX state of \cite{ChenLiuWen2011CZX} is placed on a torus of
$2N\times 2M$ qubits, blocked into $N\times M$ sites of four qubits, and is
the product over plaquettes of the states $\ket{0000}+\ket{1111}$; its tensor
is $\sum_{i,j,k,l}\ket{ijkl}((i,j),(j,k),(k,l),(l,i)|$ of bond dimension
four.\footnote{Source traceability:
    \path{Papers/2011.12127/TN-Review-main.tex}, lines 2502--2514;
    \path{References/1106.4752/source/dDSPTmodel.tex}, lines 317--321.}
The AKLT state places a singlet on every link of the square lattice and
projects the four spin-$\tfrac12$ at each site onto their symmetric
subspace, and the nearest-neighbour RVB state is the superposition of all
coverings of the lattice by nearest-neighbour singlets; the review writes
both as PEPS with tensors of the form
$A=P(\mathbb 1\otimes\mathbb 1\otimes Y\otimes Y)$, again with no
restriction on the lattice size.\footnote{Source traceability:
    \path{Papers/2011.12127/TN-Review-main.tex}, lines 2432--2448.}
A classical model with nearest-neighbour interaction $h$ at inverse
temperature $\beta$ gives the tensor
$A=[\sum_i\ket{i}(i\,i\,i\,i|](\mathbb 1\otimes\mathbb 1\otimes M\otimes M)$
with $M=\sum_{i,j}e^{-\beta h(i,j)/2}|i)(j|$, on the square lattice with no
restriction on its size.\footnote{Source traceability:
    \path{Papers/2011.12127/TN-Review-main.tex}, lines 2483--2494.}
All five constructions are local, so they apply for every $N,M\ge 1$ once the
lattice is read as a multigraph: a site on a torus of width two has two
distinct horizontal bonds to its neighbour, one on each side.

\section{The restriction}
The formal PEPS lives on a finite simple graph, and the torus is the simple
graph on $\mathbb Z_{w}\times\mathbb Z_{h}$ with cyclic nearest-neighbour
edges. For $w=2$ the right and the left neighbour of a site coincide, and so
do the two edges joining them: the four legs of a site are not four distinct
bonds. The contraction of the five tensors on this graph is therefore a
different network from the source's at $w=2$ or $h=2$. For the RVB tensor
the difference is visible on the coefficients: at $w=2$ the right and left
legs of a site carry the same bond label, so a horizontal dimer would occupy
two legs of the same site and the tensor vanishes on it. On a torus of
$2\times 3$ sites the formal network then has only vertical dimers, which
cannot cover a cycle of three sites, and its state is zero, while the
source's state, with its two horizontal bonds between the two sites of a
row, is a nonzero superposition of horizontal coverings. For
$w,h\ge 3$ the right and up edges of the sites are pairwise distinct and
exhaust the edge set, and the formal statements hold there:
\begin{enumerate}
    \item the cluster tensor on a torus of $N=wh$ sites has coefficient
    $2^{-N}(-1)^{\sum_{\langle u,v\rangle}\sigma_u\sigma_v}$, the sum running
    over nearest-neighbour bonds;
    \item the CZX tensor has coefficient $1$ at a configuration in which the
    four qubits around every plaquette agree, and $0$ otherwise;
    \item the AKLT tensor gives the symmetric projector at every site
    applied to a singlet on every edge;
    \item the RVB tensor gives the sum over nearest-neighbour dimer
    coverings of the torus of the product of the singlets on the covered
    edges;
    \item the classical-model tensor of a nearest-neighbour interaction $h$
    at inverse temperature $\beta$ has coefficient $e^{-\beta H(\sigma)/2}$,
    where $H$ sums $h$ over the nearest-neighbour bonds, so its diagonal
    expectation values are those of the Gibbs state $e^{-\beta H}/Z$.
\end{enumerate}
The GHZ example is not restricted: its coefficient depends only on whether
the configuration is constant, and it is proved on every torus with
$w,h\ge 2$.\footnote{Results: \leanid{TNLean.PEPS.stateCoeff_clusterPEPS},
    \leanid{TNLean.PEPS.stateCoeff_czxPEPS},
    \leanid{TNLean.PEPS.stateCoeff_akltPEPS},
    \leanid{TNLean.PEPS.stateCoeff_rvbPEPS},
    \leanid{TNLean.PEPS.stateCoeff_classicalPEPS},
    \leanid{TNLean.PEPS.stateCoeff_ghzPEPS}.}

\section{Elimination plan}
Formalize PEPS on finite multigraphs, where a torus of width two has two
edges between neighbouring sites, and restate the five examples there. The
proofs carry over unchanged once the edges of the torus are parametrized by
the right and up edges of the sites.

\bibliographystyle{alpha}
\bibliography{references}
\end{document}
